\documentclass{article}

% NeurIPS 2024 style. Download from https://neurips.cc/Conferences/2024/CallForPapers
% and place neurips_2024.sty alongside this file. Compile with:
%   pdflatex paper.tex
\usepackage[final]{neurips_2024}

\usepackage[utf8]{inputenc}
\usepackage[T1]{fontenc}
\usepackage{hyperref}
\usepackage{url}
\usepackage{booktabs}
\usepackage{graphicx}
\usepackage{amsmath}
\usepackage{microtype}

\title{An Inverse-MAE Ensemble for 18-Month Daily Sales Forecasting\\
  in a Vietnamese Fashion Retailer}

\author{%
  Duy \\
  DATATHON 2026 Round 1 --- Part 3: Sales Forecasting
}

\begin{document}
\maketitle

\begin{abstract}
We forecast daily Revenue and COGS for 548 days (2023-01-01 to 2024-07-01)
given a decade of transactional history (2012-07-04 to 2022-12-31). A data
audit shows that of the 14 provided tables only \texttt{sales.csv} spans the
test horizon, so the task reduces to univariate forecasting with calendar
covariates. We ensemble three contrasting models --- a seasonal-naive baseline
with compound YoY growth, a trend-detrended LightGBM on Fourier features, and
a Prophet model with log-transformed targets and anomaly-derived holidays ---
combined via inverse-MAE weights from a yearly walk-forward cross-validation.
The ensemble improves Revenue MAE by $\sim$0.4\% and COGS MAE by $\sim$0.08\%
over the best individual model on held-out 2020--2022 folds. SHAP analysis
confirms that the annual Fourier component and a 2-year Revenue lag carry the
bulk of the predictive signal.
\end{abstract}

\section{Problem and Data}

The organizer provides 14 CSVs of e-commerce activity. A data audit reveals
that only \texttt{sales.csv} contains rows after 2022-12-31, so any external
aggregate (orders, traffic, promotions, inventory) can act only on the train
side. The forecast target is the pair $(Revenue_t, COGS_t)$ for each
$t \in [2023\text{-}01\text{-}01, 2024\text{-}07\text{-}01]$.

\paragraph{Constraints.} (i) no external data, (ii) no leakage from test-period
targets, (iii) reproducibility (fixed seed 42), (iv) explainability.

\paragraph{Trend.} Annual revenue is non-monotonic: it peaks in 2016, bottoms
near 2019, and partially recovers by 2022. The geometric YoY growth across the
train span is 0.962, slightly deflationary. Any model assuming monotonic growth
misextrapolates toward 2024.

\section{Method}

\subsection{Feature engineering}
Test-time features are reconstructed from \texttt{Date} only:
\begin{itemize}
  \item Calendar: year, month, day, day-of-week, day-of-year, week, quarter,
    and weekend / month-end / quarter-end / year-end flags.
  \item Fourier terms on periods 365.25 (order 6), 7 (order 3), 30.44 (order 3).
  \item Lags at 548, 729, 1094 days and rolling mean/std (windows 7, 28, 90)
    shifted by 548 so every rolling value is drawn from pre-test history.
  \item Anomaly-day flag: a binary indicator for (month, day) pairs whose
    Revenue deviates by $|z| > 2$ from a 30-day rolling baseline in $\ge 3$
    distinct training years. This proxies Vietnamese holidays (T\'{e}t,
    30/4, 2/9, Black Friday) without external calendars.
  \item Promo-active count: number of \texttt{promotions.csv} windows live on
    each date (projects only where \texttt{end\_date} $\ge$ test start).
\end{itemize}

\subsection{Base models}
\textbf{Seasonal-naive baseline.} Let $\bar S_{md}$ be the average of the
normalized revenue (divided by its annual mean) for month $m$, day $d$ across
training years. The prediction at $t = (y, m, d)$ is
$\hat y_t = B \cdot g^{y - y_0} \cdot \bar S_{md}$,
where $B = \text{annual total}(y_0) / 365$ and $g$ is the geometric YoY
growth across full training years.

\textbf{Detrended LightGBM.} We divide each training target by
$B \cdot g^{y - y_0}$ and train LightGBM to predict the resulting seasonal
ratio (mean $\approx 1$). Year and linear-time columns are dropped to force
LGBM to learn seasonality only; the trend is reattached at inference.

\textbf{Prophet.} Log1p-transformed target with multiplicative yearly and
weekly seasonality, changepoint prior 0.05, and the anomaly-day set fed as a
single ``Holidays'' component.

\subsection{Ensemble}
Let $e_m$ be the mean CV MAE of model $m$ across folds
$\{2020, 2021, 2022\}$. Weights are
$w_m = (1/e_m) / \sum_k (1/e_k)$, and the final prediction is
$\hat y_t = \sum_m w_m \hat y_t^{(m)}$, applied independently for Revenue and
COGS.

\subsection{Cross-validation}
Walk-forward expanding-window: train on $[2012, y-1]$, validate on the full
calendar year $y$, for $y \in \{2020, 2021, 2022\}$. Final fits use all
$[2012, 2022]$ data.

\section{Experiments}

\begin{table}[h]
  \centering
  \caption{Mean CV MAE across the 2020--2022 yearly walk-forward folds.}
  \begin{tabular}{lrr}
    \toprule
    Model & Revenue MAE & COGS MAE \\
    \midrule
    Seasonal-naive baseline & 610{,}020 & 523{,}701 \\
    Detrended LightGBM      & 800{,}923 & 705{,}244 \\
    Prophet                 & 860{,}218 & 707{,}396 \\
    \midrule
    Inverse-MAE ensemble    & \textbf{607{,}514} & \textbf{523{,}273} \\
    \bottomrule
  \end{tabular}
\end{table}

The baseline dominates because its explicit trend $\times$ seasonality
decomposition matches the data-generating structure. Tree models cannot
extrapolate trend without leaking it through a column; our detrending removes
this leak but loses some bias correction. Prophet's piecewise-linear trend
overfits the 2016 peak.

\section{Explainability}

SHAP on the detrended LightGBM (gain-based, 2022 as held-out):
\texttt{Revenue\_lag729}, \texttt{y\_cos1} (annual), \texttt{day\_of\_year},
\texttt{day}, \texttt{is\_anomaly\_day} collectively explain most of the
variance --- confirming that a 2-year lag plus annual Fourier is the essential
signal. See \texttt{reports/figures/shap\_summary\_revenue.png}.

Prophet's component plot (\texttt{reports/figures/prophet\_components\_revenue.png})
shows a slightly declining trend, a sharp annual peak in Q4 (Black Friday /
Vietnamese T\'{e}t) and a pronounced weekly cycle with weekend dips. Anomaly
days contribute a multiplicative uplift of 10--25\% during holiday windows.

\section{Business insights}

\begin{itemize}
  \item \textbf{Q4 dominates the year}: annual Fourier peaks in December;
    budget fulfillment and stock planning should front-load Q4.
  \item \textbf{Weekday cycle is strong}: weekends run $\sim$15--20\% below
    weekday revenue --- staffing can be adjusted.
  \item \textbf{Trend has reset twice}: after the 2016 peak and again after
    the 2019 trough. A strict compound-growth assumption understates 2022--2024
    if the recovery continues. We report a conservative forecast anchored to
    the 10-year geometric mean.
\end{itemize}

\section{Limitations}

(i) Only one source table covers the test window, capping achievable MAE.
(ii) The 18-month horizon outruns any short lag; we compensate with calendar
structure but cannot capture macroeconomic or competitor shifts.
(iii) Ensemble gain over the baseline is small; a single-model submission
would be nearly identical in expected leaderboard score.

\section*{Reproducibility}
All code, configs and figures at
\url{https://github.com/<USER>/datathon-2026-part3}.
Running \texttt{python run\_pipeline.py} after restoring \texttt{data/}
regenerates every submission CSV and figure used in this paper.

\end{document}
